% this is the original plan for evaluation. Not included in main.tex. 

\section{Evaluation}
% Include your evaluation methodology. The philosopher of science Karl Popper states that a theory is scientific only if it is falsifiable — if some experiment exists that would disprove it. This section should include the research questions (e.g., an experimentally refutable thesis or hypothesis) and how you will answer them. You should also state your evaluation metric, which should be, at least in part, quantifiable. Another way to think of this is: how will you convince a skeptical reader that you succeeded? They aren't going to just take your word that your project was a success. It is best to include as many tables or graphs as your final report will have (with empty body for now, but with captions and with axes labeled) to show how you will communicate any quantitative results.

We will evaluate our side-effect analysis tool to answer three primary research questions regarding its scalability, correctness, and practical utility in downstream software engineering tasks.

\begin{itemize}
    \item RQ1: Scalability: What are the performance characteristics and scalability limits of the Salcianu--Rinard analysis on modern Java libraries?

    \item RQ2: Correctness: How precise and accurate is our tool in identifying side-effect-free methods compared to existing specifications and manual verification?

    \item RQ3: Downstream Utility: Does the inferred side-effect information improve the precision of downstream software engineering tools?
\end{itemize}

\textbf{RQ1}: Since Madhavan et al. \cite{MadhavanPurity2011} suggests that the algorithm described in the 2005 paper may suffer from significant performance bottlenecks, it is critical to establish an empirical baseline for its performance on modern Java bytecode. We evaluate our analysis on core JDK packages to determine the practical feasibility of this approach at the scale of standard library code. For each method we analyze, we measure  
\begin{itemize}  
\item \textit{analysis time}: the time required to run the entire analysis, 
\item \textit{call graph time}: the time required to generate the call dependency graph,
\item \textit{graph time}: the time required to construct a point-to graph at the end of the method from Jimple code, and
\item \textit{verification time}: time to determine side-effect information from point-to graph 
\end{itemize}

This evaluation will determine the practical boundary of the analysis --- identifying whether it can handle the full JDK or if it is effectively limited to smaller, application-level libraries. If the runtime is impractically long, then we might consider implementing simple optimization, including node merging, adapted from the 2011 Microsoft paper \cite{MadhavanPurity2011}.
\\

\textbf{RQ2:} We will verify the correctness of our analysis through two different ground truths:

\begin{itemize}
    
\item Comparison with JDK Specifications: We compare our inferred side-effect verdicts against existing manual purity annotations (such as \texttt{@SideEffectFree}) in the Checker Framework Annotated JDK. We report Annotation Deficit rather than standard recall because the manually annotated JDK is not an exhaustive ground truth. Since many methods are not yet manually annotated, we cannot definitively identify all false negatives; we can only identify cases where our tool finds side-effect-free methods that were previously undocumented. 

Discrepancies are categorized as either ``Tool False Positive' (the tool missed a side effect) or ``Annotation Deficit' (the tool found a side-effect-free method not yet manually annotated). Sample data table is shown below:

\begin{table}[h]
\centering
\caption{Comparison of Side-Effect Verdicts against JDK Specifications (RQ2)}
\label{tab:purity-results}
\begin{tabular}{lrr}
\hline
\textbf{Category} & \textbf{Count} & \textbf{Percentage} \\ \hline
Matches & 1,240 & 85.0\% \\
Tool False Positives & 12 & 0.8\% \\
Annotation Deficit & 208 & 14.2\% \\ \hline
\textbf{Total Methods Tested} & \textbf{1,460} & \textbf{100.0\%} \\ \hline
\end{tabular}
\end{table}
Ideally, we will show the table for each core JDK library that we test. 

\item Extensive Unit Testing: We will self-generate a comprehensive suite of unit tests covering tricky Java idioms, including aliasing, loops, and shallow/deep copy. These tests will serve as a regression suite to ensure our implementation strictly adheres to our definition of side-effect-free.
\end{itemize}

\textbf{RQ3}: We will demonstrate the utility of our tool by integrating its output into two of the three widely used software analysis tools. Ideally, if time allows, we will test all three. 
\begin{itemize}
    \item \textbf{The Checker Framework}. The Checker Framework relies on side-effect information to maintain precision. For example, if the analyzer knows a field F is \texttt{NonNull}, but then encounters a call to an unknown method \texttt{F.unknownMethod()}, it must conservatively assume \texttt{F} might have become null to maintain soundness. But precision is lost. 
    To evaluate the utility of our tool on Checker Framework, we will run the Checker Framework using: (1) current JDK side-effect annotations, and (2) current specifications plus our inferred side-effect stubs. We will then be able to compare:
\begin{itemize}
    \item Total warnings and distinct warning sites
    \item Reduction in warnings attributable to lost flow refinements across calls
    \item Any regressions (new warnings or type-check failures).
\end{itemize}

\item \textbf{Randoop (Safe Test Generation)}. Randoop is an automated unit test generator for Java that produces tests in JUnit format. To guarantee test validity, Randoop must avoid using side-effecting methods within assertions (e.g., \texttt{assert(obj.isValid())}). Invoking a side-effecting method inside an assertion alters the program state, effectively invalidating the test case it attempts to verify.

Randoop accepts a whitelist of side-effect-free methods via the \texttt{--side-effect-free-methods} argument. Providing this information enables the tool to safely utilize more methods in assertions, leading to more meaningful and stable tests. Notably, the Randoop developers explicitly list improved side-effect analysis as a key area for future enhancement \cite{RandoopProjectIdeas}.

To evaluate our tool's impact, we will execute Randoop in two configurations: \begin{itemize} \item \textbf{Run A (Baseline):} Randoop with default settings. \item \textbf{Run B (Enhanced):} Randoop configured with the \texttt{--side-effect-free-methods} list generated by our analysis. \end{itemize}

We will assess whether our analysis effectively filters out methods that have side effects, preventing Randoop from generating ``flaky'' or invalid tests that inadvertently mutate the state they are intended to verify.


\item \textbf{Daikon (Invariant Detection)}. Daikon \cite{DaikonWebsite} infers invariants (e.g., \texttt{x > 0}) by observing program execution. It needs to call observer methods to inspect object state, but calling a method that may have side effects during observation is unsafe. It could crash the program or corrupt the data being measured.

To complete the experiment and test the utility of our tool, we will run Daikon with the \texttt{--purity-file} flag, feeding it the side-effect results from our tool.

We will measure if providing this automated purity list allows Daikon to safely execute more observer methods, resulting in a richer set of detected invariants compared to running Daikon without this information.

\end{itemize}